\section{结论}
\label{sec:8}

本文围绕「算力约束下如何配置资源以提升大语言模型能力」，建立四个层层衔接的模型，主要结论如下。

\textbf{问题一。} 以 A1 冻结的 CRITIC 加权 Huber 中心作为唯一质量口径，把 22 个字段展开为 25 项指标，得到文档、领域与语料质量；六来源诊断在三集各圈定约 $5\%$ 的高分歧复核候选，主导对立方向随语料变化；17 域质量由映射与文本桥接补齐（五折 $R^2=0.4106$）。配比模型给出 $17\times13$ 迁移结构，推荐配比 $\bp^*$ 在 1M 代理模型下较均匀配比预测降损 $9.47\%$（$95\%$ 区间 $[7.82\%,11.16\%]$）；域质量与训练边际价值秩相关仅 $-0.211$，质量分不能替代配比实验。

\textbf{问题二。} 建立幂次广义标度律 $L=E+k_0(1-Q)+[AN^{-\alpha}Q^{-\kappa_N}+BD^{-\beta}Q^{-\kappa_D}]\exp(s_{\rm red}\tilde\varphi(\bp))$，其中经典参数 $E=1.6901$、$A=406.39$、$\alpha=0.34$、$B=410.59$、$\beta=0.28$，质量参数 $\kappa_N=0.1416$、$\kappa_D=0.0112$、$k_0=0.2280$、$Q_0=0.6624117$。参照点上质量弹性约为参数弹性的 $1.62$ 倍；质量提高 $0.1$ 等效于参数扩大 $1.30$–$3.07$ 倍；沿算力前沿等效算力乘 $1.191/1.485/2.134$。$C\ge10^{20}$ 后提质比扩参数更有效；配比通道受单一尺度限制，按探索性下、上情景报告。

\textbf{问题三。} 约束取等后解析消元，把问题降为 $(\log_{10}N,Q)$ 二维有界优化，以质量约束活跃状态定义结构性转移，以 $\Phi=1$ 作为固定 $Q$ 下的局部一阶必要条件（导数核验误差约 $4.1\times10^{-10}$）。上下文长度临界值为 $L_{ctx}^{\rm crit}=30000$ token。数值上，$C=10^{19}$ 时指数型、幂函数型成本下不提质、对数型直接提满，$C=10^{22},10^{24}$ 时均提满；离开下界的临界预算（$\log_{10}C$）约为 $19.06/19.53/18.58$，为「何时投向提质」给出随成本函数移动的明确阈值。

\textbf{问题四。} 用「标度律＋S 形桥接＋连续时间项」分解历史前沿：$2023.50\to2025.00$ 窗口 base 观测提高 $24.759$ 分，规模贡献 $23.063$、技术贡献 $3.404$，技术占比 $12.86\%$（$90\%$ 区间 $6.07\%$–$21.96\%$）；去饱和综合分同期提高 $48.19\%$，更灵敏地反映进展。算力增速由 $0.60$ 降至 $0.15$ dex/年时，24 个月 base 中位数为 $44.11$，技术增量占比由 $54.94\%$ 升至 $76.70\%$。该结果以 $T_0=2025.25$ 为历史锚点，更新数据后需重作。

四问共同指向核心判断：\textbf{数据质量与领域配比是与算力并列的一等资源}。低预算阶段边际投入主要投向模型规模，预算跨过质量收益与提质开销的平衡点后，提质与配比优化优先级上升；算力增长放缓时，能力进步将越来越依赖非规模技术创新。上述定量判断的适用条件已在 \ref{sec:7.2} 逐条界定。
